\subsection{局部常值层}

设 B 为拓扑空间，F 为一个集；给 F 赋以离散拓扑。称以 F 为茎型的 B 上常值层为 B 上以 F 为值的局部常值函数所成的层 (I, p. 45, example 2)。当不致担心与空间 B 相混淆时，这个层有时记作 \(\underline{\mathbf{F}}\)。它与平展 B-空间 \((B \times F,\mathrm{pr}_1)\) 的连续截面在 B 上所成的层相恒同：公式 \(i([U,s,a]) = (a,s(a))\) 确实定义了一个典范同构 \(i\) （从与 \(\underline{\mathbf{F}}\) 相伴的平展空间到 \((B \times F,\mathrm{pr}_1)\) 上的）。对每个 \(a \in B\)，映射 \([U,s,a] \mapsto s(a)\) 是茎 \(\underline{\mathbf{F}}_a\) （\(\underline{\mathbf{F}}\) 在 \(a\) 处的）到集 F 上的典范双射。

设 \(\mathcal{P} = (\mathcal{P}(\mathrm{U}), r_{\mathrm{UV}})\) 为 B 上的预层，使得对 B 的每个开集 U 有 \(\mathcal{P}(\mathrm{U}) = \mathrm{F}\)，且对 B 的每对开集 (U, V) 有 \(r_{\mathrm{UV}} = \mathrm{Id}_{\mathrm{F}}\) （只要 \(\mathrm{U} \subset \mathrm{V}\)）。于是，层 \(\widetilde{\mathcal{P}}\) （与 \(\mathcal{P}\) 相伴的）典范地同构于常值层 \(\underline{\mathrm{F}}\) (I, p. 52, example 1)。

定义 3. --- 拓扑空间 B 上的层 \(\mathcal{F}\) 是局部常值的，如果 B 的每一点都有开邻域 U，使得诱导层 \(\mathcal{F} \mid U\) 同构于 U 上的一个常值层。

命题 9. --- 一个层是局部常值的，当且仅当与之相伴的平展空间是覆叠。

设 B 为拓扑空间，\(\mathcal{F}\) 为 B 上的层。层 \(\mathcal{F}\) 是常值的，当且仅当平展空间 \(\mathrm{E}_{\mathcal{F}}\) B-同构于平凡覆叠 \((\mathrm{B} \times \mathrm{F},\mathrm{pr}_1)\)，其中 F 是离散拓扑空间。若 U 是 B 的开子集，则与诱导层 \(\mathcal{F} \mid \mathrm{U}\) 相伴的平展空间与 \(\mathrm{E}_{\mathcal{F}}\) 在 U 上诱导的 U-平展空间相恒同（参看 I, p. 51）。命题由此得出。

推论. --- 一个平展 B-空间是覆叠，当且仅当其截面的层是局部常值的。

这由该命题以及 I, p. 52 的例 2 得出。

